% Sección para el informe de la Entrega 2. Se incluye con % Sección para el informe de la Entrega 2. Se incluye con % Sección para el informe de la Entrega 2. Se incluye con % Sección para el informe de la Entrega 2. Se incluye con \input{fase2_seccion_adquisicion}
% (requiere booktabs y listings, ya cargados en main.tex). Los valores salen de docs/calidad/qc_mgn_pdet.md.
\section{Adquisición y Verificación de los Datos}

\subsection{Fuentes y mecanismo de descarga}
Los límites municipales se toman del Marco Geoestadístico Nacional 2025 del DANE. En el plan de la Entrega~1 se
previó descargar el ZIP completo (1.5~GB, todos los niveles geográficos), pero el geoportal publica el nivel
municipio por separado (\texttt{MGN2025\_MPIO\_GRAFICO.zip}, 71.6~MB, Shapefile en WGS84) con la misma versión y la
misma capa, por lo que se descarga solo ese archivo. El listado de los 170 municipios PDET y su subregión se toma
de la Agencia de Renovación del Territorio (Decreto 893 de 2017).

La descarga la hace \texttt{src/00a\_download\_mgn.py}: lee el archivo por bloques de 1~MB a un temporal
\texttt{.part}, reintenta y reanuda con cabecera \texttt{Range} si la conexión se corta, y solo renombra el
archivo cuando su tamaño coincide con el \texttt{Content-Length} del servidor. La URL, las rutas y los valores de
referencia están en \texttt{config/acquisition.yaml}, no en el código.

\begin{lstlisting}[style=lab]
python src/00a_download_mgn.py      # descarga y verifica (si ya existe, solo verifica)
python src/00c_filter_pdet.py       # filtro PDET y controles -> data/mgn_pdet_2025.gpkg
python src/01_load_municipalities.py
\end{lstlisting}

\subsection{Verificación de integridad}
El DANE no publica un \emph{checksum} oficial. Por eso, en la primera descarga (08/11/2026) se fijó en la
configuración el SHA-256 del ZIP
(\texttt{a3c01393\ldots bc141f}), su tamaño (71\,576\,494 bytes) y las cabeceras \texttt{Last-Modified} y
\texttt{ETag}. Cada corrida recalcula el SHA-256 y falla si no coincide, lo que garantiza que todo el equipo
procese exactamente los mismos bytes y detecta si el DANE reemplaza el archivo. Además se valida el CRC de los ocho
archivos del ZIP y la presencia de los componentes del Shapefile (\texttt{.shp}, \texttt{.shx}, \texttt{.dbf},
\texttt{.prj}, \texttt{.cpg}). El SHA-256 de cada archivo extraído queda en \texttt{docs/calidad/manifest\_mgn.json}
y se vuelve a comprobar antes del filtrado.

\subsection{Filtro PDET y controles de calidad}
\texttt{src/00c\_filter\_pdet.py} cruza la capa municipal (1\,122 unidades) con el listado PDET por código
DIVIPOLA. En el Excel de la ART los códigos de Antioquia aparecen sin el cero inicial (\texttt{5490} en lugar de
\texttt{05490}), así que se normalizan a cinco dígitos antes del cruce. El Cuadro~\ref{tab:qc_mgn} resume los
32 controles; el detalle está en \texttt{docs/calidad/qc\_mgn\_pdet.md}.

\begin{table}[!htbp]
\centering\small
\begin{tabular}{@{}p{2.3cm}p{10.2cm}c@{}}
\toprule
Grupo & Controles & Resultado \\
\midrule
Fuente & SHA-256 del ZIP verificado; archivos extraídos sin cambios & 2 OK \\
Capa MGN & Campos requeridos; CRS EPSG:4326; 1\,122 unidades; código único, de 5 dígitos y coherente con el
departamento; sin nulos; extensión dentro de Colombia & 8 OK \\
Listado PDET & 170 municipios; códigos únicos y de 5 dígitos; 16 subregiones; prefijo coherente; sin nulos & 6 OK \\
Cruce & 170/170 códigos encontrados; departamento coincide; nombre coincide & 2 OK, 1 AVISO \\
Geometría & 170 polígonos únicos, no vacíos y válidos; área EPSG:9377 vs.\ \texttt{mpio\_narea} (diferencia
0.000\,\%); sin solapes (339 pares vecinos); geometría más pesada 2.2~MB (límite BSON 16~MB) & 10 OK \\
Salida & GeoPackage releído: 170 registros, EPSG:4326, sin nulos & 3 OK \\
\bottomrule
\end{tabular}
\caption{Controles de calidad sobre el MGN2025 filtrado a los municipios PDET (31 OK, 1 AVISO, 0 FALLA).}
\label{tab:qc_mgn}
\end{table}

El único aviso corresponde a seis municipios cuyo nombre en el listado de la ART difiere del nombre oficial del
MGN2025 (por ejemplo, \emph{Vista Hermosa} frente a \emph{Vistahermosa} y \emph{Tolú Viejo} frente a \emph{San
José de Toluviejo}). Como el cruce es por código, no afecta el resultado; se conserva el nombre del MGN. Los 170
municipios suman 389\,182~km$^2$ en 16 subregiones. La salida, \texttt{data/mgn\_pdet\_2025.gpkg} (EPSG:4326,
geometrías \texttt{MultiPolygon}), es la entrada del ETL de integridad y se probó con él en modo
\texttt{--dry-run}: 170 municipios listos, sin geometrías rechazadas.

% (requiere booktabs y listings, ya cargados en main.tex). Los valores salen de docs/calidad/qc_mgn_pdet.md.
\section{Adquisición y Verificación de los Datos}

\subsection{Fuentes y mecanismo de descarga}
Los límites municipales se toman del Marco Geoestadístico Nacional 2025 del DANE. En el plan de la Entrega~1 se
previó descargar el ZIP completo (1.5~GB, todos los niveles geográficos), pero el geoportal publica el nivel
municipio por separado (\texttt{MGN2025\_MPIO\_GRAFICO.zip}, 71.6~MB, Shapefile en WGS84) con la misma versión y la
misma capa, por lo que se descarga solo ese archivo. El listado de los 170 municipios PDET y su subregión se toma
de la Agencia de Renovación del Territorio (Decreto 893 de 2017).

La descarga la hace \texttt{src/00a\_download\_mgn.py}: lee el archivo por bloques de 1~MB a un temporal
\texttt{.part}, reintenta y reanuda con cabecera \texttt{Range} si la conexión se corta, y solo renombra el
archivo cuando su tamaño coincide con el \texttt{Content-Length} del servidor. La URL, las rutas y los valores de
referencia están en \texttt{config/acquisition.yaml}, no en el código.

\begin{lstlisting}[style=lab]
python src/00a_download_mgn.py      # descarga y verifica (si ya existe, solo verifica)
python src/00c_filter_pdet.py       # filtro PDET y controles -> data/mgn_pdet_2025.gpkg
python src/01_load_municipalities.py
\end{lstlisting}

\subsection{Verificación de integridad}
El DANE no publica un \emph{checksum} oficial. Por eso, en la primera descarga (08/11/2026) se fijó en la
configuración el SHA-256 del ZIP
(\texttt{a3c01393\ldots bc141f}), su tamaño (71\,576\,494 bytes) y las cabeceras \texttt{Last-Modified} y
\texttt{ETag}. Cada corrida recalcula el SHA-256 y falla si no coincide, lo que garantiza que todo el equipo
procese exactamente los mismos bytes y detecta si el DANE reemplaza el archivo. Además se valida el CRC de los ocho
archivos del ZIP y la presencia de los componentes del Shapefile (\texttt{.shp}, \texttt{.shx}, \texttt{.dbf},
\texttt{.prj}, \texttt{.cpg}). El SHA-256 de cada archivo extraído queda en \texttt{docs/calidad/manifest\_mgn.json}
y se vuelve a comprobar antes del filtrado.

\subsection{Filtro PDET y controles de calidad}
\texttt{src/00c\_filter\_pdet.py} cruza la capa municipal (1\,122 unidades) con el listado PDET por código
DIVIPOLA. En el Excel de la ART los códigos de Antioquia aparecen sin el cero inicial (\texttt{5490} en lugar de
\texttt{05490}), así que se normalizan a cinco dígitos antes del cruce. El Cuadro~\ref{tab:qc_mgn} resume los
32 controles; el detalle está en \texttt{docs/calidad/qc\_mgn\_pdet.md}.

\begin{table}[!htbp]
\centering\small
\begin{tabular}{@{}p{2.3cm}p{10.2cm}c@{}}
\toprule
Grupo & Controles & Resultado \\
\midrule
Fuente & SHA-256 del ZIP verificado; archivos extraídos sin cambios & 2 OK \\
Capa MGN & Campos requeridos; CRS EPSG:4326; 1\,122 unidades; código único, de 5 dígitos y coherente con el
departamento; sin nulos; extensión dentro de Colombia & 8 OK \\
Listado PDET & 170 municipios; códigos únicos y de 5 dígitos; 16 subregiones; prefijo coherente; sin nulos & 6 OK \\
Cruce & 170/170 códigos encontrados; departamento coincide; nombre coincide & 2 OK, 1 AVISO \\
Geometría & 170 polígonos únicos, no vacíos y válidos; área EPSG:9377 vs.\ \texttt{mpio\_narea} (diferencia
0.000\,\%); sin solapes (339 pares vecinos); geometría más pesada 2.2~MB (límite BSON 16~MB) & 10 OK \\
Salida & GeoPackage releído: 170 registros, EPSG:4326, sin nulos & 3 OK \\
\bottomrule
\end{tabular}
\caption{Controles de calidad sobre el MGN2025 filtrado a los municipios PDET (31 OK, 1 AVISO, 0 FALLA).}
\label{tab:qc_mgn}
\end{table}

El único aviso corresponde a seis municipios cuyo nombre en el listado de la ART difiere del nombre oficial del
MGN2025 (por ejemplo, \emph{Vista Hermosa} frente a \emph{Vistahermosa} y \emph{Tolú Viejo} frente a \emph{San
José de Toluviejo}). Como el cruce es por código, no afecta el resultado; se conserva el nombre del MGN. Los 170
municipios suman 389\,182~km$^2$ en 16 subregiones. La salida, \texttt{data/mgn\_pdet\_2025.gpkg} (EPSG:4326,
geometrías \texttt{MultiPolygon}), es la entrada del ETL de integridad y se probó con él en modo
\texttt{--dry-run}: 170 municipios listos, sin geometrías rechazadas.

% (requiere booktabs y listings, ya cargados en main.tex). Los valores salen de docs/calidad/qc_mgn_pdet.md.
\section{Adquisición y Verificación de los Datos}

\subsection{Fuentes y mecanismo de descarga}
Los límites municipales se toman del Marco Geoestadístico Nacional 2025 del DANE. En el plan de la Entrega~1 se
previó descargar el ZIP completo (1.5~GB, todos los niveles geográficos), pero el geoportal publica el nivel
municipio por separado (\texttt{MGN2025\_MPIO\_GRAFICO.zip}, 71.6~MB, Shapefile en WGS84) con la misma versión y la
misma capa, por lo que se descarga solo ese archivo. El listado de los 170 municipios PDET y su subregión se toma
de la Agencia de Renovación del Territorio (Decreto 893 de 2017).

La descarga la hace \texttt{src/00a\_download\_mgn.py}: lee el archivo por bloques de 1~MB a un temporal
\texttt{.part}, reintenta y reanuda con cabecera \texttt{Range} si la conexión se corta, y solo renombra el
archivo cuando su tamaño coincide con el \texttt{Content-Length} del servidor. La URL, las rutas y los valores de
referencia están en \texttt{config/acquisition.yaml}, no en el código.

\begin{lstlisting}[style=lab]
python src/00a_download_mgn.py      # descarga y verifica (si ya existe, solo verifica)
python src/00c_filter_pdet.py       # filtro PDET y controles -> data/mgn_pdet_2025.gpkg
python src/01_load_municipalities.py
\end{lstlisting}

\subsection{Verificación de integridad}
El DANE no publica un \emph{checksum} oficial. Por eso, en la primera descarga (08/11/2026) se fijó en la
configuración el SHA-256 del ZIP
(\texttt{a3c01393\ldots bc141f}), su tamaño (71\,576\,494 bytes) y las cabeceras \texttt{Last-Modified} y
\texttt{ETag}. Cada corrida recalcula el SHA-256 y falla si no coincide, lo que garantiza que todo el equipo
procese exactamente los mismos bytes y detecta si el DANE reemplaza el archivo. Además se valida el CRC de los ocho
archivos del ZIP y la presencia de los componentes del Shapefile (\texttt{.shp}, \texttt{.shx}, \texttt{.dbf},
\texttt{.prj}, \texttt{.cpg}). El SHA-256 de cada archivo extraído queda en \texttt{docs/calidad/manifest\_mgn.json}
y se vuelve a comprobar antes del filtrado.

\subsection{Filtro PDET y controles de calidad}
\texttt{src/00c\_filter\_pdet.py} cruza la capa municipal (1\,122 unidades) con el listado PDET por código
DIVIPOLA. En el Excel de la ART los códigos de Antioquia aparecen sin el cero inicial (\texttt{5490} en lugar de
\texttt{05490}), así que se normalizan a cinco dígitos antes del cruce. El Cuadro~\ref{tab:qc_mgn} resume los
32 controles; el detalle está en \texttt{docs/calidad/qc\_mgn\_pdet.md}.

\begin{table}[!htbp]
\centering\small
\begin{tabular}{@{}p{2.3cm}p{10.2cm}c@{}}
\toprule
Grupo & Controles & Resultado \\
\midrule
Fuente & SHA-256 del ZIP verificado; archivos extraídos sin cambios & 2 OK \\
Capa MGN & Campos requeridos; CRS EPSG:4326; 1\,122 unidades; código único, de 5 dígitos y coherente con el
departamento; sin nulos; extensión dentro de Colombia & 8 OK \\
Listado PDET & 170 municipios; códigos únicos y de 5 dígitos; 16 subregiones; prefijo coherente; sin nulos & 6 OK \\
Cruce & 170/170 códigos encontrados; departamento coincide; nombre coincide & 2 OK, 1 AVISO \\
Geometría & 170 polígonos únicos, no vacíos y válidos; área EPSG:9377 vs.\ \texttt{mpio\_narea} (diferencia
0.000\,\%); sin solapes (339 pares vecinos); geometría más pesada 2.2~MB (límite BSON 16~MB) & 10 OK \\
Salida & GeoPackage releído: 170 registros, EPSG:4326, sin nulos & 3 OK \\
\bottomrule
\end{tabular}
\caption{Controles de calidad sobre el MGN2025 filtrado a los municipios PDET (31 OK, 1 AVISO, 0 FALLA).}
\label{tab:qc_mgn}
\end{table}

El único aviso corresponde a seis municipios cuyo nombre en el listado de la ART difiere del nombre oficial del
MGN2025 (por ejemplo, \emph{Vista Hermosa} frente a \emph{Vistahermosa} y \emph{Tolú Viejo} frente a \emph{San
José de Toluviejo}). Como el cruce es por código, no afecta el resultado; se conserva el nombre del MGN. Los 170
municipios suman 389\,182~km$^2$ en 16 subregiones. La salida, \texttt{data/mgn\_pdet\_2025.gpkg} (EPSG:4326,
geometrías \texttt{MultiPolygon}), es la entrada del ETL de integridad y se probó con él en modo
\texttt{--dry-run}: 170 municipios listos, sin geometrías rechazadas.

% (requiere booktabs y listings, ya cargados en main.tex). Los valores salen de docs/calidad/qc_mgn_pdet.md.
\section{Adquisición y Verificación de los Datos}

\subsection{Fuentes y mecanismo de descarga}
Los límites municipales se toman del Marco Geoestadístico Nacional 2025 del DANE. En el plan de la Entrega~1 se
previó descargar el ZIP completo (1.5~GB, todos los niveles geográficos), pero el geoportal publica el nivel
municipio por separado (\texttt{MGN2025\_MPIO\_GRAFICO.zip}, 71.6~MB, Shapefile en WGS84) con la misma versión y la
misma capa, por lo que se descarga solo ese archivo. El listado de los 170 municipios PDET y su subregión se toma
de la Agencia de Renovación del Territorio (Decreto 893 de 2017).

La descarga la hace \texttt{src/00a\_download\_mgn.py}: lee el archivo por bloques de 1~MB a un temporal
\texttt{.part}, reintenta y reanuda con cabecera \texttt{Range} si la conexión se corta, y solo renombra el
archivo cuando su tamaño coincide con el \texttt{Content-Length} del servidor. La URL, las rutas y los valores de
referencia están en \texttt{config/acquisition.yaml}, no en el código.

\begin{lstlisting}[style=lab]
python src/00a_download_mgn.py      # descarga y verifica (si ya existe, solo verifica)
python src/00c_filter_pdet.py       # filtro PDET y controles -> data/mgn_pdet_2025.gpkg
python src/01_load_municipalities.py
\end{lstlisting}

\subsection{Verificación de integridad}
El DANE no publica un \emph{checksum} oficial. Por eso, en la primera descarga (08/11/2026) se fijó en la
configuración el SHA-256 del ZIP
(\texttt{a3c01393\ldots bc141f}), su tamaño (71\,576\,494 bytes) y las cabeceras \texttt{Last-Modified} y
\texttt{ETag}. Cada corrida recalcula el SHA-256 y falla si no coincide, lo que garantiza que todo el equipo
procese exactamente los mismos bytes y detecta si el DANE reemplaza el archivo. Además se valida el CRC de los ocho
archivos del ZIP y la presencia de los componentes del Shapefile (\texttt{.shp}, \texttt{.shx}, \texttt{.dbf},
\texttt{.prj}, \texttt{.cpg}). El SHA-256 de cada archivo extraído queda en \texttt{docs/calidad/manifest\_mgn.json}
y se vuelve a comprobar antes del filtrado.

\subsection{Filtro PDET y controles de calidad}
\texttt{src/00c\_filter\_pdet.py} cruza la capa municipal (1\,122 unidades) con el listado PDET por código
DIVIPOLA. En el Excel de la ART los códigos de Antioquia aparecen sin el cero inicial (\texttt{5490} en lugar de
\texttt{05490}), así que se normalizan a cinco dígitos antes del cruce. El Cuadro~\ref{tab:qc_mgn} resume los
32 controles; el detalle está en \texttt{docs/calidad/qc\_mgn\_pdet.md}.

\begin{table}[!htbp]
\centering\small
\begin{tabular}{@{}p{2.3cm}p{10.2cm}c@{}}
\toprule
Grupo & Controles & Resultado \\
\midrule
Fuente & SHA-256 del ZIP verificado; archivos extraídos sin cambios & 2 OK \\
Capa MGN & Campos requeridos; CRS EPSG:4326; 1\,122 unidades; código único, de 5 dígitos y coherente con el
departamento; sin nulos; extensión dentro de Colombia & 8 OK \\
Listado PDET & 170 municipios; códigos únicos y de 5 dígitos; 16 subregiones; prefijo coherente; sin nulos & 6 OK \\
Cruce & 170/170 códigos encontrados; departamento coincide; nombre coincide & 2 OK, 1 AVISO \\
Geometría & 170 polígonos únicos, no vacíos y válidos; área EPSG:9377 vs.\ \texttt{mpio\_narea} (diferencia
0.000\,\%); sin solapes (339 pares vecinos); geometría más pesada 2.2~MB (límite BSON 16~MB) & 10 OK \\
Salida & GeoPackage releído: 170 registros, EPSG:4326, sin nulos & 3 OK \\
\bottomrule
\end{tabular}
\caption{Controles de calidad sobre el MGN2025 filtrado a los municipios PDET (31 OK, 1 AVISO, 0 FALLA).}
\label{tab:qc_mgn}
\end{table}

El único aviso corresponde a seis municipios cuyo nombre en el listado de la ART difiere del nombre oficial del
MGN2025 (por ejemplo, \emph{Vista Hermosa} frente a \emph{Vistahermosa} y \emph{Tolú Viejo} frente a \emph{San
José de Toluviejo}). Como el cruce es por código, no afecta el resultado; se conserva el nombre del MGN. Los 170
municipios suman 389\,182~km$^2$ en 16 subregiones. La salida, \texttt{data/mgn\_pdet\_2025.gpkg} (EPSG:4326,
geometrías \texttt{MultiPolygon}), es la entrada del ETL de integridad y se probó con él en modo
\texttt{--dry-run}: 170 municipios listos, sin geometrías rechazadas.
